Athletics (track and field) carries a substantial burden of health problems at the competitive tiers of participation \citep{mckay2022tier}. A systematic review and meta-analysis of 216 studies in Olympic athletics disciplines reported a mean injured-athlete prevalence of 50.4\% (range 2.4\%--93.6\%) and a mean incidence of 135.6 injuries per 1000 registered athletes, with injury patterns differing markedly by discipline group, tissue type, and period of the season \citep{edouard2020injury}. Roughly half of injuries in youth athletics academies were reported in the winter preparation period \citep{martinezsilvan2021injury}, motivating separate description of preparation- and competition-period injury rather than a single season-level estimate.

How injury is registered materially changes the burden observed. Surveillance restricted to absence, complete sports cessation, or medical attention can under-ascertain overuse problems because athletes may continue training and competing despite symptoms and impaired function \citep{clarsen2013ostrc}. The Oslo Sports Trauma Research Center Questionnaire on Health Problems (OSTRC-H) therefore complements time-loss registration with direct athlete reports across four impact domains: sports participation, training volume, performance, and symptoms \citep{clarsen2014ostrch,clarsen2020ostrc}. These domains provide an athlete-reported measure of sport-specific health burden; they should not be interpreted as a generic quality-of-life scale.

Since 2020--2021, competitive athletics has undergone a rapid equipment transition with the introduction of ``advanced footwear technology'' (AFT) spikes: a rigid, longitudinally curved carbon-fiber plate embedded in a high-resilience foam midsole, first validated in road racing (Nike Vaporfly, Breaking2 project) and subsequently extended to track spikes for sprint and middle-distance events \citep{healey2022superspikes}. The stated biomechanical rationale is that plate-driven longitudinal bending stiffness reduces energy dissipation at the metatarsophalangeal joints and shifts the center of pressure anteriorly, increasing horizontal propulsive force \citep{willwacher2016stiffness}. Comparative analyses of world-leading sprint performance before and after the introduction of AFT spikes have documented step changes in top-100 performance lists coincident with the technology's diffusion \citep{mason2023superspikes}.

The same mechanical changes that plausibly improve performance also alter loading at the ankle and forefoot. A recent kinetic comparison of traditional and carbon-plated sprint spikes in elite sprinters found materially higher ankle joint contact force and higher peak plantarflexor muscle force in the carbon-plated condition, with pronounced inter-individual variability \citep{gerlach2025aft}. Case reports and narrative reviews have raised the possibility of increased Achilles tendon load and of altered forces at the navicular and cuneiform bones, both recognized high-risk sites for bone stress injury in runners \citep{tenforde2023bonestress}. Evidence on ground-reaction-force-mediated risk is mixed, and no consensus exists on whether increased longitudinal bending stiffness is protective, neutral, or hazardous once inter-individual differences in plantarflexor strength are taken into account \citep{willwacher2016stiffness,gerlach2025aft}.

Despite this plausible mechanistic pathway, direct epidemiological evidence on carbon-plate insole use and injury in competitive athletics is essentially absent: available data are limited to isolated case reports and small mechanistic cohorts, none of which were designed to estimate population-level injury associations \citep{tenforde2023bonestress,gerlach2025aft}. Established risk factors for running-related injury more broadly include age, body mass index, weekly training frequency, footwear, years of sporting experience, and prior injury \citep{taunton2003running}, all of which needed to be accounted for before any exposure--outcome association attributable to carbon-plate use could be considered informative.

This study was designed to close part of this gap by describing carbon-plate insole use among a national sample of Italian competitive athletics athletes and by estimating the cross-sectional, period-matched association between reported carbon share and self-reported injury during the preparation and competition periods of a single season. Consistent with the Injury Definitions Concept Framework and the athletics-specific consensus statement on injury and illness surveillance \citep{timpka2014consensus,timpka2015meta}, injury was defined broadly as a physical complaint or observable tissue damage attributable to the transfer of energy experienced by an athlete during training or competition participation in athletics, irrespective of whether medical attention was sought or of the consequences for participation \citep{bahr2020strobesiis}. The primary hypothesis, formulated a priori, was that greater reported carbon share would be associated with higher adjusted injury prevalence; a secondary, exploratory aim was to describe whether carbon-plate adoption and share differed by age group and sex, and whether the exposure--outcome association was stable under population reweighting and alternative dose--response specifications.
